本文从谱范数约束下的更新方向出发，完整重构了五步经验NS、动量与权重衰减下的有限步驻点界。
证明通过方向对齐而非单调正交化误差控制下降，保留初始化、稳定项与参数安排的作用。
正式实验完成45次矩阵优化和9次小型语言模型训练。
固定步长的矩阵问题中，AdamW在所列网格下终点目标值更低；语言模型中，NS5的平均测试损失为1.5626，AdamW为1.6073。
这种跨任务差异要求结合问题结构和参数安排理解优化器，不能由一个任务推断普遍优劣。
进一步扩大数据和模型规模、增加种子及调参预算，才有条件研究结论的稳定性。
